\section{Two steps toward the sparse fourth moment}
\label{app:iwanieclocal}

This section proves the passage to a central divisor polynomial and
its first mean-square localization. These are two steps in the
nonspectral approach of Jutila~\cite{Jut91}. The arithmetic mean-square
estimate is the additional cited input identified at the end of this
section. The moment results use \cite[(\ref*{six-eq:basefourth})]{Weighted} as stated.

\begin{proposition}[central divisor shell by contour shifting]
\label{prop:centralshell}
Fix $a=1/(8\pi)$ and $b=4/\pi$, and put
\[
 D_T(v)=\sum_{aT<n\le bT}\frac{d(n)}{n^{1/2+iv}},\qquad H=\log T.
\]
Uniformly for $T\le t\le2T$, with $T$ sufficiently large,
\begin{equation}\label{eq:centralshell}
 |\zeta(1/2+it)|^4
 \ll (\log T)^2\left(1+H\int_{|u|\le H^2}
                    e^{-u^2/H^2}|D_T(t+u)|^2\,du\right).
\end{equation}
All endpoints of the shell and all its coefficients are fixed before
$t$ varies.
\end{proposition}
The proof follows directly from the functional equation.
\begin{proof}
Write the proved functional equation as
$\zeta(s)=\chi(s)\zeta(1-s)$, where
\[
 \chi(s)^{-1}=\pi^{1/2-s}
             \frac{\Gamma(s/2)}{\Gamma((1-s)/2)}.
\]
On the critical line, $|\chi(s)|=1$ and
$\chi(s)^{-1}\zeta(s)^2=|\zeta(s)|^2$.
We first prove, for every fixed $A>0$,
\begin{equation}\label{eq:shellGaussian}
 \int_{\mathbb R}e^{-(v-t)^2/H^2}|\zeta(1/2+iv)|^2\,dv
 =\int_{\mathbb R}e^{-(v-t)^2/H^2}
       \chi(1/2+iv)^{-1}D_T(v)\,dv+O_A(T^{-A}).
\end{equation}

Insert $W(s)=\exp((s-1/2-it)^2/H^2)$ and move its vertical
integral from $\Re s=1/2$ to $\Re s=2$.
The only pole in the strip of $\chi(s)^{-1}\zeta(s)^2$ is at $s=1$.
It is simple, since $\chi(s)^{-1}=-(s-1)/2+O((s-1)^2)$ there.
Its weighted residue is $O(\exp((1/4-t^2)/H^2))$.
Euler summation and the gamma bounds give polynomial growth on this
fixed strip, so horizontal sides vanish. On $\Re s=2$ the absolutely
convergent series $\zeta(s)^2=\sum d(n)n^{-s}$ may be integrated
term by term. Denote its $n$th kernel by
\[
 K_t(n)=\frac1i\int_{\Re s=2}W(s)\chi(s)^{-1}n^{-s}\,ds.
\]
We show
\begin{equation}\label{eq:shelltail}
 \sum_{n\notin(aT,bT]}d(n)|K_t(n)|\ll_A T^{-A}.
\end{equation}
This justifies discarding the complete complement of the shell.
The finite shell integrand is holomorphic across the contour shift,
proving \eqref{eq:shellGaussian}.

Here are uniform details of the kernel estimate. Put $s_0=1/2+it$,
$q=t/(2\pi)$, and choose $\theta(t)\in\mathbb R$ with
$\chi(s_0)^{-1}=e^{i\theta(t)}$. In the disk $|z|\le c\sqrt T$,
with $c>0$ fixed and small, the function
\[
 F_t(z)=e^{-i\theta(t)}\chi(s_0+z)^{-1}q^{-z}
\]
is holomorphic, has $F_t(0)=1$, and is bounded uniformly in $t,T$.
For this bound, differentiate
the Euler-summation formula for $\log\Gamma$ used in
\cite[Appendix~\ref*{six-app:localzeta}]{Weighted}. On this disk its first derivative gives
\[
 (\log\chi^{-1})'(s_0)=\log q+O(T^{-1}),\qquad
 (\log\chi^{-1})''(s_0+z)=O(T^{-1}).
\]
Both gamma arguments stay a distance $\gg T$ from their poles and
inside fixed sectors away from the negative real axis. The Euler remainder formula continues analytically to the plane cut
along the negative real axis. Its differentiated remainder integrals
are bounded here using $|x+z|^2\gg T^2+(x+\Re z)^2$,
for $x\ge0$ and either gamma argument $z$. This proves the estimates
also when the real part of an argument is negative. Taylor's integral formula now gives
$\log F_t(z)=O(|z|/T+|z|^2/T)$.
Cauchy's formula on this disk gives, for any fixed integer $B$,
\[
 F_t(z)=\sum_{j=0}^{B}c_j(t)z^j+
 O_B\bigl(T^{-(B+1)/2}(1+|z|)^{B+1}\bigr),
 \qquad c_j(t)\ll_j T^{-j/2},
\]
for $|z|\le H^2+2$.
On the line $\Re z=3/2$ the part $|\Im z|>H^2$ costs
$O_A(T^{-A}n^{-2})$, by the Gaussian and polynomial gamma growth.
For the Taylor remainder the cost is
$O_B(n^{-2}T^{3/2-(B+1)/2}H^{B+2})$.
Taking $B$ large in terms of $A$ bounds its sum against $d(n)$.
Each polynomial term is an elementary differentiated Gaussian integral:
\[
 K_t(n)=\frac{e^{i\theta(t)-it\log n}}{\sqrt n}
 \sum_{j=0}^{B}c_j(t)\,\partial_\lambda^j
       \{\sqrt\pi H e^{-H^2\lambda^2/4}\}
       +O_A(T^{-A}n^{-2}),
 \qquad \lambda=\log(q/n).
\]
Here the last $A$ can be chosen arbitrarily by increasing $B$.
For $n\notin(aT,bT]$ we have $|\lambda|\ge\log4$.
The differentiated Gaussian is at most
$C_B H^{2B+1}(1+|\lambda|)^B e^{-H^2\lambda^2/4}$.
Grouping integers into intervals with $|\log(n/q)|$ between
successive fixed positive multiples, and using
$\sum_{n\le X}d(n)\ll X\log(2X)$, bounds the leading tail by
$T^{1/2}H^{O_B(1)}e^{-c_0H^2}$.
This is $O_A(T^{-A})$ for every fixed $A$, and proves
\eqref{eq:shelltail}, including the entire large-$n$ tail.

\cite[Proposition~\ref*{six-prop:localzetaproof}]{Weighted}, with $p=2$, bounds the local
value by $O(\log T)$ times one plus the left side of
\eqref{eq:shellGaussian}; the Gaussian is bounded below on the
local interval. Take absolute values on its right, use $|\chi|=1$,
and square. Cauchy's inequality and
$\int e^{-u^2/H^2}du=\sqrt\pi H$ prove \eqref{eq:centralshell}
with the integral over the whole line. Finally
$|D_T(v)|\ll\sqrt T\log T$ for all real $v$, by the divisor identity.
The portion $|u|>H^2$ is smaller than any negative power of $T$,
which proves the stated version.
\end{proof}

\begin{lemma}[exact short-frequency localization]
\label{lem:dirichletGallagher}
Let $b_n$ be any finite complex sequence, let $L\ge1$, let $t_0$ be
real, and set $h=1/(2L)$. Then
\begin{equation}\label{eq:shortfrequency}
 \int_{t_0-L/2}^{t_0+L/2}
       \left|\sum_n b_n n^{-it}\right|^2dt
 \ll L^2\int_0^\infty
       \left|\sum_{x<n\le e^h x}b_n n^{-it_0}\right|^2\frac{dx}{x}.
\end{equation}
\end{lemma}
The inequality keeps every endpoint and every coefficient $b_n$.
\begin{proof}
Put
\[
 A(y)=\sum_{y<\log n\le y+h}b_n n^{-it_0}.
\]
Use the radian Fourier transform $\widehat A(u)=\int A(y)e^{-iuy}dy$.
Direct integration of each interval gives exactly
\[
 \widehat A(u)=e^{iuh/2}\frac{2\sin(uh/2)}u
                  \sum_n b_n n^{-i(t_0+u)}.
\]
For $|u|\le L/2$ the absolute value of the first multiplier is
at least a fixed positive constant times $h$.
Plancherel therefore proves the inequality, followed by $x=e^y$.
To prove this form of Plancherel, first convolve
the finite step function $A$ with a Gaussian. Fourier inversion for
the resulting smooth function, followed by Fubini, gives the norm
identity. Letting the Gaussian width tend to zero gives the step
function identity by convergence in $L^2$. This proves the stated Plancherel step directly.
\end{proof}

\subsection{The arithmetic mean-square estimate}
The preceding two proofs identify the precise arithmetic input.
The bound to be supplied by the cited mean-square theorem is,
for disjoint intervals of
length $L$ in a fixed dilation of $[T,2T]$,
\begin{equation}\label{eq:dirichletmeansquare}
 \sum_{r\le R}\int_{I_r}|D_T(v)|^2\,dv
 \prec RL+T\sqrt{R/L},\qquad \sqrt T\le L\le T.
\end{equation}
Indeed, integrate \eqref{eq:centralshell} over the intervals $I_r$.
Their $H^2$-enlargements have bounded overlap since $L\ge\sqrt T$;
the Gaussian integration costs only logarithms. Applying
\eqref{eq:dirichletmeansquare} to boundedly many disjoint classes yields
\cite[(\ref*{six-eq:basefourth})]{Weighted} with the same exponents.
For $L$ near $\sqrt T$ one may first cover the intervals by grid cells
of length $T^\delta L$, with at most $2R$ occupied cells. This costs
$T^\delta$ and places the length in Jutila's interior range.
For $L\ge T^{2/3}$, subdividing into lengths comparable to $T^{2/3}$
gives the required bound $O(RL)$. These two operations explain the
usual $T^\varepsilon$ endpoint losses while keeping every interval.

In Jutila's Theorem~1, take a single additive twist, $M\asymp T$,
$f(x,v,y)=-(y+v)\log x/(2\pi)$, and $g(x)=x^{-1/2}$.
The factor $(G+MG')^2M$ is bounded, $F\asymp T$, $V=L$, and the
second term of (1.13) becomes $T\sqrt{R/L}$; finitely many dyadic
pieces cover the fixed shell. His proof of that mean-square estimate
uses the additive-divisor transformation (Lemma~4), the two simultaneous
spacing conditions for fractions and their modular inverses (Lemma~3),
and the double large sieve (Lemma~5). Lemma~3 itself refers to the
iteration and reflection arguments of Huxley--Watt and to its
predecessor in Jutila's earlier paper. Those cited arguments give
the arithmetic off-diagonal estimate in \eqref{eq:dirichletmeansquare}.
This section gives \eqref{eq:shortfrequency} and
\eqref{eq:centralshell}; either formulation keeps the stated
mean-square theorem as its analytic input.
